%---------------------------------------------------------------
\chapter{Implementace}
%---------------------------------------------------------------

\section{Postup při realizaci projektu}

\subsection{Příprava implementace}

\subsubsection{Aktualizace konvencí, linterů}

Na začátku přípravy implementace jsem se zaměřil na zavedení nových konvencí a linterů s cílem zvýšit kvalitu kódu a zajistit jeho konzistenci napříč celým projektem. V projektu byly již zavedeny knihovny ESLint a Prettier, nicméně jejich konfigurace byla zastaralá a neodpovídala současným standardům. Proto jsem aktualizoval obě knihovny, jejich konfigurace a zavedl jsem nové konvence a pravidla pro psaní kódu. Po zavedení těchto změn jsem provedl refaktoring stávajícího kódu, aby odpovídal novým pravidlům. 

\subsubsection{Dockerizace}

Následně jsem provedl dockerizaci vývojového prostředí, která do té doby nebyla v aplikaci zavedena. Cílem bylo zajistit, aby veškerý vývoj probíhal v konzistentním prostředí a aby se zjednodušilo nastavování aplikace a práce s jednotlivými příkazy, skripty a knihovnami.

\subsubsection{Zavední nástrojů umělé inteligence}

% todo

\subsubsection{Design System}

Po zavedení nových konvencí a linterů jsem se přesunul k vytvoření design systému.